\section{Introduction}
\label{sec:introduction}
Repository-level software agents increasingly combine code search, editing, command execution, and feedback within one workflow. SWE-bench measures the ability of complete systems to resolve real repository issues, while SWE-agent and OpenHands demonstrate the importance of the interfaces and execution environments surrounding a language model \cite{jimenez2024swebench,yang2024sweagent,wang2025openhands}. These systems make a second question consequential: when an agent acts on a previously observed result, what connects that result to the software state being changed?

Consider an agent modifying session-refresh behavior. It locates the correct function and obtains its source. Another worker then updates the same file. A test result remains available from the earlier revision, a second verifier reports a pass for a narrower target, and the implementation task is marked complete. None of these observations is necessarily false. The error arises if their distinct meanings are collapsed into a single indication that the current change is ready. Localization does not establish that an edit's source is current; an executed check does not establish coverage of every requirement; task completion does not establish acceptance.

This paper studies that coordination problem at the runtime boundary. We use \emph{evidence} to mean a typed observation about an artifact or operation, with recorded provenance and conditions of applicability. Evidence may remain useful as history after it ceases to justify a current operation. The key distinction is therefore between retaining an observation and admitting it into a present decision. This distinction matters whenever an agent, a human editor, and background checks share evolving software state.

We present \wcode{}, an implemented Rust runtime exposed through the Model Context Protocol. It maintains declared software intent, code observations, verification records, and work state. Its organizing rule is that an operation consumes evidence with explicit applicability conditions. Reading returns original source and file identity; mutation checks the current file against the observed identity; verification records the revision and authority under which a check ran; reconciliation preserves relevant disagreements between producers. The same rule also determines what the interface reports as missing or uncertain.

The research question is: \emph{How can a shared software-agent runtime preserve the applicability of evidence across bounded observation, concurrent modification, and verification?} We examine it through three contributions:
\begin{enumerate}
\item An explicit admission contract relating evidence to code--design revision, target scope, authority, and policy. The contract distinguishes current support, historical observations, and unresolved producer conflicts.
\item A runtime implementation that carries these distinctions across context packing, guarded single-file edits, verification execution and reuse, and reconciliation. A running example shows why preserving the fields at each boundary is necessary.
\item An auditable diagnostic evaluation separating source delivery, edit-input readiness, operational controls, and descriptive cost. A retained executable and source archive identify the primary run; a predicate-removal analysis exposes a readiness error hidden by saturated retrieval scores.
\end{enumerate}

The contribution is a composition of established mechanisms around a shared operational contract. Hashing, provenance, architecture conformance, and access control each have substantial prior work; Section~\ref{sec:related} locates the runtime within those traditions. Our evaluation is correspondingly specific. It invokes no language model and establishes behavior on developer-authored scenarios rather than improved patch success. This scope allows us to inspect failures at tool boundaries without attributing them to model reasoning, while leaving downstream effectiveness to a separate experiment.
